% Accéder au cluster Zen : compte MesoNET, clé SSH, première connexion
% Source : documentation MesoNET (mesonet.fr), consultée le 16/11/2026.
\documentclass[8pt,aspectratio=169]{beamer}

\usepackage[french]{babel}
\usepackage[T1]{fontenc}
\usepackage{microtype}
\usepackage{booktabs,array}
\usepackage{tikz}
\usetikzlibrary{babel,arrows.meta}

\usetheme[progressbar=frametitle,numbering=fraction,block=fill]{metropolis}
\metroset{sectionpage=none,titleformat=regular}
\setbeamersize{text margin left=0.55cm,text margin right=0.55cm}

\definecolor{hpcblue}{HTML}{1F4E79}
\definecolor{hpcorange}{HTML}{C55A11}
\definecolor{hpcgreen}{HTML}{3B7D23}
\definecolor{hpcgrey}{HTML}{D9D9D9}
\setbeamercolor{palette primary}{bg=hpcblue,fg=white}
\setbeamercolor{progress bar}{fg=hpcorange,bg=hpcgrey}
\setbeamercolor{frametitle}{bg=hpcblue,fg=white}

\newcolumntype{L}[1]{>{\raggedright\arraybackslash}p{#1}}
\newcommand{\src}[1]{\vfill{\tiny\color{gray}Source : #1\par}}
\setlength{\parskip}{2pt}
\setbeamertemplate{itemize items}[circle]

\title{Accéder au cluster Zen}
\subtitle{Compte MesoNET, clé SSH, première connexion}
\author{HPC --- 2026-2027}
\date{Séance 1}

\begin{document}

\maketitle

% ---------------------------------------------------------------
\begin{frame}{Objectif, vue d'ensemble et prérequis}
\begin{columns}[T]
\begin{column}{0.55\textwidth}
À la fin de cette partie, vous avez :
\begin{enumerate}
\item un \textbf{compte MesoNET} rattaché à votre établissement (eduGAIN) ;
\item un \textbf{profil} complété dans le projet \emph{Zen -- Formation Zen} ;
\item une \textbf{paire de clés SSH}, clé publique déposée sur le portail ;
\item un \textbf{premier \texttt{ssh}} réussi vers \texttt{zen.univ-lille.fr}.
\end{enumerate}

\bigskip
\centering
\resizebox{\linewidth}{!}{%
\begin{tikzpicture}[
  wf/.style={draw=hpcblue,fill=hpcblue!12,rounded corners=3pt,minimum width=1.85cm,minimum height=1.0cm,align=center,font=\footnotesize},
  arr/.style={-{Stealth},thick,hpcorange}]
 \node[wf] (a) at (0,0) {1. Mail\\ d'invitation};
 \node[wf] (b) at (2.2,0) {2. Compte\\ MesoNET};
 \node[wf] (c) at (4.4,0) {3. Profil\\ Zen};
 \node[wf] (d) at (6.6,0) {4. Clés\\ SSH};
 \node[wf,fill=hpcgreen!20] (e) at (8.8,0) {5. Connexion\\ à Zen};
 \draw[arr] (a)--(b); \draw[arr] (b)--(c); \draw[arr] (c)--(d); \draw[arr] (d)--(e);
\end{tikzpicture}}

\raggedright\smallskip
Étapes 1 à 3 dans le navigateur, étape 4 dans un terminal puis sur le portail, étape 5 dans un terminal.
\end{column}
\begin{column}{0.43\textwidth}
\textbf{Prérequis}
\begin{itemize}
\item Être inscrit à la formation : \textbf{le responsable vous a invité} sur MesoNET.
\item Avoir accès à l'adresse mail renseignée par le responsable (réception de l'invitation).
\item Un terminal avec un client SSH :
 \begin{itemize}
 \item \textbf{Linux / macOS} : déjà installé ;
 \item \textbf{Windows} : MobaXterm (Home, version portable).
 \end{itemize}
\end{itemize}
\begin{block}{Note}
Aucun mot de passe sur Zen : connexion \textbf{obligatoirement par clé SSH}.
\end{block}
\begin{alertblock}{Avertissement}
Le mail d'invitation a une \textbf{date d'expiration}. Une invitation expirée ne permet plus l'accès : prévenez-moi rapidement.
\end{alertblock}
\end{column}
\end{columns}
\src{documentation MesoNET, pages \og accès étudiants \fg{} et \og connexion \fg{} de Zen (liens en dernière diapo).}
\end{frame}

% ---------------------------------------------------------------
\begin{frame}{Étapes 1 à 3 : mail, compte MesoNET, profil}
\begin{columns}[T]
\begin{column}{0.5\textwidth}
\textbf{1. Le mail d'invitation}

\smallskip
Deux liens, à utiliser \textbf{dans l'ordre} :

\smallskip
\begin{tabular}{@{}cL{2.6cm}L{3.3cm}@{}}
\toprule
Ordre & Lien & Usage\\
\midrule
1 & Création de compte & Créer le compte MesoNET\\
2 & Profil & À ouvrir \textbf{après} validation du compte\\
\bottomrule
\end{tabular}

\bigskip
\textbf{2. Créer le compte MesoNET}

\smallskip
Cliquez sur le \textbf{premier lien} et suivez la procédure :
\begin{itemize}
\item \textbf{Identification eduGAIN} : tapez quelques lettres, votre établissement apparaît dans le menu.
 \begin{itemize}
 \item Absent\,? Demandez à votre responsable. \textbf{Pas de compte local.}
 \end{itemize}
\item \textbf{Champ \og Notes \fg} (obligatoire) : \textbf{intitulé de la formation} et \textbf{établissement}.
\end{itemize}
\end{column}
\begin{column}{0.46\textwidth}
\textbf{3. Compléter le profil}

\smallskip
Après validation du compte, ouvrez le \textbf{second lien} :

\smallskip
\begin{tabular}{@{}L{2.5cm}L{4.2cm}@{}}
\toprule
Champ & Valeur\\
\midrule
Laboratoire & \texttt{Zen - Formation Zen}\\
Statut & \texttt{étudiant}\\
Établissement & votre établissement\\
\bottomrule
\end{tabular}

\smallskip
En cas de doute sur l'établissement, demandez à votre responsable de formation.

\medskip
\begin{block}{Exemple de champ \og Notes \fg}
\emph{Master \dots{} --- cours HPC --- Université Côte d'Azur}. Sans ce champ, la demande ne peut pas être traitée.
\end{block}
\begin{exampleblock}{Astuce}
Ne cliquez pas sur le second lien tant que le compte n'est pas validé.
\end{exampleblock}
\end{column}
\end{columns}
\src{documentation MesoNET, \og accès étudiants \fg{} de Zen.}
\end{frame}

% ---------------------------------------------------------------
\begin{frame}[fragile]{Étape 4a : générer sa paire de clés SSH}
\begin{columns}[T]
\begin{column}{0.5\textwidth}
\textbf{Principe} : clé \textbf{privée} (personnelle, ne la partagez jamais, elle reste sur votre machine), clé \textbf{publique} \texttt{.pub} (déposée sur le portail), \textbf{passphrase} pour protéger la clé privée. Type recommandé : \textbf{ed25519}.

\medskip
\textbf{Linux / macOS}
\begin{verbatim}
ssh-keygen -t ed25519 -f ~/.ssh/id_ed25519
cat ~/.ssh/id_ed25519.pub
\end{verbatim}
\begin{itemize}
\item Choisissez une \textbf{passphrase} quand elle est demandée.
\item \texttt{id\_ed25519} = clé \textbf{privée}, \texttt{id\_ed25519.pub} = clé \textbf{publique}.
\end{itemize}

Si ed25519 ne fonctionne pas, variante RSA (\textbf{3072 bits minimum}) :
\begin{verbatim}
ssh-keygen -t rsa -b 3072 -f ~/.ssh/id_rsa
\end{verbatim}
\end{column}
\begin{column}{0.46\textwidth}
\textbf{Windows (MobaXterm)}
\begin{enumerate}
\item Téléchargez la version portable sur \texttt{mobaxterm.mobatek.net}, extrayez le \texttt{.zip}.
\item Lancez MobaXterm, ouvrez \textbf{MobaKeyGen}.
\item Choisissez \textbf{EdDSA (Ed25519)}.
\item \textbf{Generate}, bougez la souris jusqu'à ce que la barre verte soit pleine.
\item \textbf{Key comment} explicite (par ex. \texttt{Zen}) et \textbf{passphrase}.
\item \textbf{Save private key} (conservez-la\,!), puis \textbf{Save public key}.
\item \textbf{Laissez la fenêtre ouverte} : vous copierez le cadre \og clé publique \fg{} dans le portail.
\end{enumerate}
\end{column}
\end{columns}
\src{documentation MesoNET, \og clés SSH \fg.}
\end{frame}

% ---------------------------------------------------------------
\begin{frame}{Étape 4b : déposer la clé et l'associer au projet}
\begin{columns}[T]
\begin{column}{0.5\textwidth}
\textbf{Déposer la clé publique}
\begin{enumerate}
\item \textbf{acces.mesonet.fr} $\to$ \textbf{Connexion} (via eduGAIN, comme à la création du compte).
\item Icône \textbf{Utilisateur} en haut à droite $\to$ \textbf{Clés ssh}.
\item \textbf{Ajouter une clé}.
\item Renseignez :
 \begin{itemize}
 \item \textbf{Nom de la clé} : un identifiant unique ;
 \item \textbf{Votre clé publique ssh} : contenu du fichier \texttt{.pub}.
 \end{itemize}
\item Cliquez sur \textbf{ajouter}.
\end{enumerate}
\end{column}
\begin{column}{0.46\textwidth}
\textbf{Associer la clé au projet}

\smallskip
Déposer la clé ne suffit pas : elle doit être \textbf{associée à votre compte dans le projet}.
\begin{enumerate}
\item Dans la \textbf{liste des projets} : \textbf{Ajouter une clé SSH}.
\item \textbf{Acceptez les CGU} de la machine.
\item \textbf{Sélectionnez la clé}, puis \textbf{Modifier}.
\end{enumerate}
\begin{block}{Note}
Une clé en attente de déploiement est signalée par \textbf{\textcolor{hpcorange}{!}}. Le symbole disparaît une fois la clé déployée ; le délai dépend de la machine.
\end{block}
\end{column}
\end{columns}
\src{documentation MesoNET, \og clés SSH \fg.}
\end{frame}

% ---------------------------------------------------------------
\begin{frame}[fragile]{Étape 5 : première connexion}
\begin{columns}[T]
\begin{column}{0.5\textwidth}
\textbf{Votre identifiant}
\begin{itemize}
\item Comptes créés \textbf{automatiquement toutes les 10 minutes} après la procédure d'accès.
\item Format \textbf{\texttt{m2XXXX-YY}} : \texttt{m2XXXX} = numéro du \textbf{projet}, \texttt{YY} = numéro d'utilisateur.
\item Un identifiant par couple \textbf{projet / machine}.
\item Le portail (liste des projets) affiche votre login et la clé associée.
\end{itemize}

\medskip
\textbf{Se connecter}
\begin{verbatim}
ssh -X m2XXXX-YY@zen.univ-lille.fr \
    -i ~/.ssh/id_ed25519
\end{verbatim}
\texttt{-i} peut être omis si la clé a un nom standard (\texttt{id\_ed25519}, \texttt{id\_rsa}). Aucun VPN n'est mentionné dans la documentation.
\end{column}
\begin{column}{0.46\textwidth}
\textbf{Raccourci avec \texttt{\textasciitilde/.ssh/config}}
\begin{verbatim}
Host zen
    Hostname zen.univ-lille.fr
    User m2XXXX-YY
    IdentityFile ~/.ssh/id_ed25519
\end{verbatim}
Puis simplement :
\begin{verbatim}
ssh zen
\end{verbatim}
\begin{block}{Note}
Sous Windows, MobaXterm : session SSH vers \texttt{zen.univ-lille.fr}, avec votre identifiant et votre clé privée.
\end{block}
\end{column}
\end{columns}
\src{documentation MesoNET, \og connexion \fg{} à Zen.}
\end{frame}

% ---------------------------------------------------------------
\begin{frame}[fragile]{Dépannage}
\begin{columns}[T]
\begin{column}{0.5\textwidth}
\begin{alertblock}{Attention}
\textbf{5 échecs de connexion} $\to$ adresse IP \textbf{bloquée 1 heure}. Après un premier échec, \textbf{attendez 10 minutes} avant de réessayer.
\end{alertblock}

Vérifiez dans l'ordre :
\begin{enumerate}
\item Le compte a-t-il eu le temps d'être créé (10 minutes)\,?
\item La clé est-elle \textbf{associée au projet} et sans symbole d'attente\,?
\item Bonne \textbf{clé privée} (\texttt{-i}) et bon \textbf{identifiant}\,?
\item Avez-vous déposé la clé \textbf{publique} (\texttt{.pub}) et non la privée\,?
\end{enumerate}
\end{column}
\begin{column}{0.46\textwidth}
\textbf{Toujours bloqué\,?}

\smallskip
Écrivez à \textbf{support@mesonet.fr} avec :
\begin{itemize}
\item votre nom et votre projet ;
\item votre adresse IP publique ;
\item la sortie de :
\end{itemize}
\begin{verbatim}
ssh -vvv -X m2XXXX-YY@zen.univ-lille.fr
\end{verbatim}
\end{column}
\end{columns}
\src{documentation MesoNET, \og connexion \fg{} à Zen.}
\end{frame}

% ---------------------------------------------------------------
\begin{frame}{Bilan et ressources}
\begin{columns}[T]
\begin{column}{0.5\textwidth}
\textbf{Checklist}
\begin{itemize}
\item[$\square$] Compte MesoNET créé via eduGAIN, champ \og Notes \fg{} rempli
\item[$\square$] Profil complété (\emph{Zen -- Formation Zen}, étudiant)
\item[$\square$] Clé SSH générée, protégée par une passphrase
\item[$\square$] Clé publique déposée sur \texttt{acces.mesonet.fr}
\item[$\square$] Clé associée au projet, CGU acceptées
\item[$\square$] \texttt{ssh} vers \texttt{zen.univ-lille.fr} réussi
\end{itemize}
\end{column}
\begin{column}{0.46\textwidth}
\textbf{Ressources}
\begin{itemize}
\item {\footnotesize\url{https://www.mesonet.fr/documentation/user-documentation/code_form/zen/acces/acces_etu}}
\item {\footnotesize\url{https://www.mesonet.fr/documentation/user-documentation/acces/ssh}}
\item {\footnotesize\url{https://www.mesonet.fr/documentation/user-documentation/code_form/zen/connexion}}
\item Portail : \texttt{acces.mesonet.fr}
\end{itemize}
\medskip
\textbf{Prochaine étape} : jobs et mini-tuto Slurm.
\end{column}
\end{columns}
\end{frame}

\end{document}
